% ECONOMICS_PROBLEM: {"id": "F13-228", "title": "Cross-border data governance", "jel_code": "F13", "source_id": "OP-0347"}
% !TeX program = xelatex
\documentclass[11pt,a4paper]{article}
\usepackage[margin=23mm]{geometry}
\usepackage{fontspec}
\setmainfont{Latin Modern Roman}
\usepackage{amsmath,amssymb,mathtools}
\usepackage{xcolor,enumitem,booktabs,longtable,array,xurl,hyperref}
\definecolor{muted}{HTML}{53636C}
\hypersetup{colorlinks=true,urlcolor=blue,linkcolor=blue}
\setlength{\parindent}{0pt}
\setlength{\parskip}{5pt}
\newcommand{\R}{\mathbb R}
\newcommand{\E}{\mathbb E}
\newcommand{\N}{\mathbb N}
\newcommand{\ind}{\mathbf 1}
\newcommand{\Var}{\operatorname{Var}}
\newcommand{\Cov}{\operatorname{Cov}}
\newcommand{\supp}{\operatorname{supp}}
\DeclareMathOperator*{\argmin}{arg\,min}
\DeclareMathOperator*{\argmax}{arg\,max}
\newcommand{\entrymeta}[3]{{\sffamily\small\color{muted}\textbf{\texttt{#1}}\quad|\quad #2\par #3\par}}
\newcommand{\Problemref}[1]{\texttt{#1}}
\newcommand{\JELref}[2]{\texttt{#2} (JEL \texttt{#1})}
\begin{document}
\section*{Cross-border data governance}
\textbf{Problem ID:} \texttt{F13-228}\quad \textbf{JEL:} \texttt{F13}\par
% BEGIN ECONOMICS STATEMENT
\label{problem:OP-0347}
{\sffamily\small\color{muted}Original subject group: Trade, globalization and geoeconomics\par}
\label{definitions:problem:30007629}
\textbf{Audit classification:} Background; proposed extension. The official report and DOI are real and its abstract gives conditional scenarios. The exact optimal cross-border governance problem is not verified as a source-stated open question; privacy/cyber valuation choices are editorial.
\subsection*{Definitions and precise research target}
Fix jurisdictions trading digitally delivered services and transmitting personal or business data. A governance regime \(a\) specifies transfer eligibility, localization, privacy safeguards, cybersecurity requirements and enforcement. Firms and users choose participation, security effort and service use in response. Let \(T(a)\) be trade-related surplus, \(P(a)\) expected privacy loss, \(S(a)\) cyber incident loss and \(C(a)\) compliance resource cost, each defined for a declared population and under a stated probability law. The task is to characterize feasible regimes maximizing
\[W(a)=T(a)-P(a)-S(a)-C(a),\qquad a\in A.\]
The set \(A\) includes sovereignty, legal and enforceability constraints, and valuations of privacy losses must be explicitly specified rather than inferred from trade volumes. Identify how safeguards affect trust and participation separately from direct changes in transaction costs. Model cross-border externalities and regulatory interaction; unilateral optimization need not yield the jointly best regime. Estimate or bound responses using policy variation with credible controls for concurrent digital development. Report welfare sensitivity to risk valuations and heterogeneous users. Complete liberalization or localization cannot be declared optimal merely by omitting trust, privacy or cybersecurity channels from the objective.

\textbf{Additional formal definitions and scope (editorial).} Personal data, nonpersonal business data and legally sensitive data are separate declared classes. A regime specifies permissible uses, storage locations, transfer safeguards, breach probability and enforcement sanctions. Let \(T\) be consumption-equivalent surplus net of ordinary production costs, \(P\) valuation of privacy harm, \(S\) valuation of cyber losses and \(C\) compliance resource costs, all in common units with nonoverlapping coverage. If cybersecurity losses already include privacy harm, subtract the overlap. Trust is a behavioral participation response, not a fifth welfare benefit to add if its consumption gains are already in \(T\). The feasible set explicitly incorporates privacy rights or hard risk caps where they cannot be traded off by monetary valuation. International regulatory interaction requires a declared cooperative benchmark or strategic regime; a scenario comparison does not identify a globally optimal law.

For the identification part, declare an admissible class \(\Theta\) of structural models, including preferences, technologies, shock laws, measurement/sampling rules and any equilibrium selector. If \(O\) denotes the complete observable history and \(T(\theta)\) the target defined above, the identified set is \(\mathcal I_T=\{T(\theta):\theta\in\Theta,\ \mathcal L_\theta(O)=\mathcal L_{\rm obs}(O)\}\); point identification means this set is a singleton. A \(do(a)\) comparison replaces stated policy/technology equations while fixing the declared exogenous law and re-solving the remaining equations. These definitions do not assert that an identification theorem holds without further restrictions.
\subsection*{Variants and assumption changes}
\begin{enumerate}
\item \textbf{Scenario trade model (source-motivated).} Compare localization and safeguarded flows with fixed estimated cost/trust responses; outcomes remain conditional on those calibrations.
\item \textbf{Privacy-constrained welfare design (editorial).} Add explicit harm valuations or nontradeable privacy constraints, cybersecurity and user heterogeneity; characterize feasible regimes and identified welfare bounds.
\end{enumerate}
\subsection*{References and source audit}
\begin{enumerate}
\item Official OECD citation verifies OECD/WTO, February 10, 2025, title and DOI. Locator: Official abstract and citation metadata. Support: Conditional data-regulation scenarios. Landing-page summary inspected; full report not verified. URL: \url{https://www.oecd.org/en/publications/economic-implications-of-data-regulation_aa285504-en.html}.
\end{enumerate}
\begin{enumerate}\raggedright
\item\hypertarget{definitions:source:30007629:1}{} OECD; World Trade Organization. 2025. \emph{Economic Implications of Data Regulation: Balancing Openness and Trust}. Official abstract and citation metadata.
\url{https://www.oecd.org/en/publications/economic-implications-of-data-regulation_aa285504-en.html}.
DOI: \url{https://doi.org/10.1787/aa285504-en}.
\end{enumerate}

\subsection*{Common conventions: Source mathematical conventions}

These conventions apply to each entry unless explicitly replaced.

\textbf{Models and identification.} A primitive model \(m\) specifies all probability laws, feasible choices, information, equilibrium and measurement rules used by its target. The admissible class \(\mathcal M\) is fixed independently of outcomes. If \(P_O(m)\) is its observable probability law and \(\tau(m)\) the target, its identified set at observable law \(P\) is
\[
 \mathcal I_\tau(P)=\{\tau(m):m\in\mathcal M,\ P_O(m)=P\}.
\]
Point identification means that this set is a singleton. An empty set signals incompatibility of data and assumptions. Coordinate bounds need not describe the joint set. Bounds are sharp only if every retained target is attainable or arbitrarily approachable by compatible models. Finite bounds require explicit restrictions on unobserved quantities; unrestricted confounding, hidden wealth, extrapolation or arbitrary equilibrium selection can yield uninformative sets.

\textbf{Causal interventions.} A potential outcome \(Y_i(p)\) is the outcome under a complete policy regime \(p\), including timing, financing, information and market spillovers. The target population and units are fixed before treatment. With interference, use the complete assignment vector or a justified exposure mapping; individual no-interference notation alone cannot identify an economy-wide intervention. A difference of mean potential outcomes does not require the cross-world joint distribution. Conditioning on post-treatment migration, survival, enrollment or employment changes the estimand and needs separate justification.

\textbf{Statistical statements.} An estimator, confidence set, forecast or test specifies a sampling law, model class and loss/coverage criterion. Uniform \((1-\alpha)\)-coverage means \(\inf_{m\in\mathcal M}\Pr_m\{\tau(m)\in C_n\}\geq1-\alpha\), or an explicitly asymptotic version. The whole parameter space trivially has coverage, so informativeness requires a stated width, risk or power objective. A forecasting improvement is relative to a named benchmark, information set and evaluation population.

\textbf{Equilibrium and optimization.} An equilibrium comprises feasible plans, optimality, beliefs where needed, budget constraints and all market/resource clearing. The primitives or a stated benchmark define each correspondence \(\mathcal E(p;m)\). Nonemptiness is assumed or proved, never inferred just from compact policy sets. A selected-equilibrium target uses a declared selection rule; a robust target quantifies over all permitted equilibria. Use suprema or infima unless attainment is established. An \(\varepsilon\)-optimal policy is feasible and within \(\varepsilon\) of the finite optimum value. Report an empty feasible set separately.

\textbf{Welfare and accounting.} Normative utility, interpersonal normalization, social weights, numeraire, discounting and the use of public receipts are primitives. Behavioral utility need not coincide with the welfare criterion. Household welfare includes ownership income and rents once; adding profits again to compensating variation generally double-counts them. A monetary equivalent is defined by an explicit transfer and price convention, with monotonicity and range conditions. Average effects, mechanism contributions, transfers and welfare losses are different targets. Infinite sums require convergence and dynamic feasible plans require terminal or no-Ponzi conditions.

\textbf{Source versions.} A working-paper date, revision date and journal publication year are distinguished. A DOI for a journal version is not automatically the DOI of an earlier manuscript. Locators refer to the stated version and printed pages where specified. A metadata-only check does not verify a claim about a full-text conclusion. Every entry's source subsection narrows the provenance claim accordingly.
% END ECONOMICS STATEMENT
\end{document}
